% ================================================================
% CHAPTER 9: PERSUASION AS PROJECTION
% ================================================================
\chapter{Persuasion as Projection}
\label{ch:persuasion}

\epigraph{The ideas of the ruling class are in every epoch the ruling ideas.}
{Karl Marx, \textit{The German Ideology}}

\section{The Geometry of Influence}

Persuasion is the act of changing someone's mind --- of moving their idea-vector
in meaning space so that it more closely aligns with the persuader's own. In the
Hilbert space model, persuasion has a strikingly clean geometric interpretation:
it is \emph{orthogonal projection}.

When a sender sends signal $s$ to a receiver with prior $r$, the receiver's
updated belief is $r' = r \compose s$, which in the Hilbert space is
$\embed{r'} = \embed{r} + \embed{s}$. The ``persuasive effect'' of $s$ on $r$
is the change in $r$'s belief: $\embed{r'} - \embed{r} = \embed{s}$. But not
all of this change is ``effective'' persuasion. The component of $\embed{s}$
parallel to $\embed{r}$ merely reinforces the receiver's existing belief; the
component orthogonal to $\embed{r}$ introduces genuinely new content.

\begin{definition}[Persuasion Decomposition]
\label{def:persuasion-decomposition}
The \textbf{persuasion decomposition} of signal $s$ relative to receiver prior
$r$ (with $\embed{r} \neq 0$) is:
\[
  \embed{s} = \underbrace{\frac{\rs{s}{r}}{\rs{r}{r}} \embed{r}}_{
  \text{reinforcement component}} +
  \underbrace{\left(\embed{s} - \frac{\rs{s}{r}}{\rs{r}{r}}
  \embed{r}\right)}_{\text{novelty component}}.
\]
The \textbf{reinforcement component} $\proj{\hat{r}}{s}$ is the orthogonal
projection of $\embed{s}$ onto the direction of $\embed{r}$. The
\textbf{novelty component} is the remainder.
\end{definition}

\begin{theorem}[Persuasion Components]
\label{thm:persuasion-components}
Let $\alpha_\parallel := \frac{\rs{s}{r}}{\sqrt{\rs{r}{r}}}$ (the signed length
of the reinforcement component) and $\alpha_\perp := \sqrt{\rs{s}{s} -
\frac{\rs{s}{r}^2}{\rs{r}{r}}}$ (the length of the novelty component). Then:
\begin{enumerate}
\item $\nrm{\embed{s}}^2 = \alpha_\parallel^2 + \alpha_\perp^2$ (Pythagorean
decomposition of signal weight);
\item The reinforcement component has norm $|\alpha_\parallel|$;
\item The novelty component has norm $\alpha_\perp$;
\item $\alpha_\parallel > 0$ iff $\rs{s}{r} > 0$ (the signal reinforces the
prior);
\item $\alpha_\parallel < 0$ iff $\rs{s}{r} < 0$ (the signal undermines the
prior).
\end{enumerate}
\end{theorem}

\begin{proof}
This is the standard decomposition theorem for orthogonal projections. The
Pythagorean identity follows because the reinforcement and novelty components
are orthogonal:
\[
  \nrm{\embed{s}}^2 = \nrm{\alpha_\parallel \hat{r} + v_\perp}^2
  = \alpha_\parallel^2 + \nrm{v_\perp}^2 = \alpha_\parallel^2 + \alpha_\perp^2.
\]

In Lean~4:
\begin{lstlisting}
theorem persuasion_pythagoras (hr : φ r ≠ 0) :
    ‖φ s‖ ^ 2 = (rs s r / ‖φ r‖) ^ 2 +
    ‖φ s - (rs s r / rs r r) • φ r‖ ^ 2 := by
  have := norm_sq_eq_inner (φ s)
  simp [rs, proj_orthogonal_sq]
  ring
\end{lstlisting}
\end{proof}

\begin{interpretation}[The Dual Nature of Persuasion]
\label{interp:dual-persuasion}
The persuasion decomposition reveals that every act of persuasion has two
distinct components:

\begin{enumerate}
\item \textbf{Reinforcement}: the part of the signal that ``preaches to the
choir'' --- aligning with what the receiver already believes. This component
increases the receiver's conviction in their existing direction without changing
that direction. It is ``effective'' in the sense that the receiver readily
absorbs it, but ``ineffective'' in the sense that it produces no genuine change
of mind.

\item \textbf{Novelty}: the part of the signal that introduces genuinely new
content --- ideas that the receiver did not previously hold. This component
changes the \emph{direction} of the receiver's belief vector (not just its
magnitude) and thus constitutes genuine persuasion. But it is ``foreign'' to
the receiver and may be resisted precisely because it does not resonate with
their prior.
\end{enumerate}

The tension between reinforcement and novelty is the central paradox of
persuasion. Effective persuaders must balance the two: enough reinforcement
to be credible and palatable to the receiver, but enough novelty to actually
change their mind. Too much reinforcement (all $\alpha_\parallel$, no
$\alpha_\perp$) is mere flattery; too much novelty (all $\alpha_\perp$, no
$\alpha_\parallel$) is incomprehensible. The optimal balance depends on the
receiver's ``openness to novelty'' --- a parameter that the basic model does
not capture but that extensions could incorporate.

In the rhetoric of Aristotle, the reinforcement component corresponds to
\emph{ethos} (the speaker's credibility, established by aligning with the
audience's values) and the novelty component corresponds to \emph{logos} (the
logical content of the argument). A skilled orator establishes ethos before
deploying logos --- in geometric terms, they project onto the audience's
direction before introducing orthogonal content.
\end{interpretation}


\section{The Effect of Persuasion on the Receiver}

\begin{theorem}[Post-Persuasion Direction Change]
\label{thm:direction-change}
After receiving signal $s$, the receiver's belief direction changes from
$\hat{r} = \embed{r}/\nrm{\embed{r}}$ to:
\[
  \hat{r}' = \frac{\embed{r} + \embed{s}}{\nrm{\embed{r} + \embed{s}}}.
\]
The angular change $\Delta\theta := \cosangle{r}{r \compose s}$ satisfies:
\[
  \cos(\Delta\theta) = \frac{\rs{r}{r} + \rs{r}{s}}{w(r) \cdot w(r \compose s)}
  = \frac{w(r) + \alpha_\parallel}{\sqrt{(w(r) + \alpha_\parallel)^2 +
  \alpha_\perp^2}},
\]
where $w(r) = \nrm{\embed{r}}$.
\end{theorem}

\begin{proof}
\begin{align*}
  \cos(\Delta\theta)
  &= \frac{\ip{\embed{r}}{\embed{r} + \embed{s}}}
          {\nrm{\embed{r}} \cdot \nrm{\embed{r} + \embed{s}}} \\
  &= \frac{\nrm{\embed{r}}^2 + \ip{\embed{r}}{\embed{s}}}
          {\nrm{\embed{r}} \cdot \nrm{\embed{r} + \embed{s}}} \\
  &= \frac{w(r)^2 + w(r) \alpha_\parallel}
          {w(r) \cdot \nrm{\embed{r} + \embed{s}}} \\
  &= \frac{w(r) + \alpha_\parallel}{\nrm{\embed{r} + \embed{s}}}.
\end{align*}
Using $\nrm{\embed{r} + \embed{s}}^2 = (w(r) + \alpha_\parallel)^2 +
\alpha_\perp^2$ (since the reinforcement component adds to $w(r)$ along the
$\hat{r}$ direction, and the novelty component is perpendicular), we get the
stated formula.
\end{proof}

\begin{corollary}[Small Novelty Approximation]
\label{cor:small-novelty}
When the novelty component is small relative to the reinforced prior
($\alpha_\perp \ll w(r) + \alpha_\parallel$):
\[
  \Delta\theta \approx \frac{\alpha_\perp}{w(r) + \alpha_\parallel}.
\]
The angular change is proportional to the novelty and inversely proportional
to the reinforced prior.
\end{corollary}

\begin{proof}
For small $x$, $\arccos(1/(1 + x^2)^{1/2}) \approx x$. Set $x =
\alpha_\perp / (w(r) + \alpha_\parallel)$.
\end{proof}

\begin{interpretation}[The Persuasion Paradox]
\label{interp:persuasion-paradox}
Corollary~\ref{cor:small-novelty} reveals a fundamental paradox of persuasion:
\emph{the stronger the receiver's prior, the harder they are to persuade}. The
angular change $\Delta\theta \approx \alpha_\perp / (w(r) + \alpha_\parallel)$
is inversely proportional to $w(r)$ --- the weight of the receiver's prior.
People with strong convictions (large $w(r)$) are hard to budge; people with
weak convictions (small $w(r)$) are easy to move.

This is geometrically intuitive: a short vector is easily deflected by adding
another vector; a long vector barely changes direction when the same vector is
added. The ``momentum'' of a long vector overwhelms the ``push'' of the
persuasive signal.

But the paradox deepens. The formula also shows that \emph{reinforcement makes
the receiver harder to persuade in the future}. The term $\alpha_\parallel$
appears in the denominator: the more the current signal reinforces the receiver's
prior, the less any subsequent novelty will change their direction. Repeated
reinforcement --- echo chambers, propaganda, ideological indoctrination ---
makes the receiver's prior stronger and stronger, making them increasingly
resistant to any future novelty.

This is a \emph{positive feedback loop}: reinforcement increases resistance to
novelty, which makes the receiver more likely to seek out reinforcement (since
novel ideas have smaller impact and thus lower payoff), which further increases
resistance. The geometric structure of the Hilbert space naturally produces
the kind of ``ideological entrenchment'' observed in polarized societies.

In Gramsci's framework, this positive feedback loop is the mechanism of
\emph{hegemony}: the dominant ideology reinforces itself through repeated
transmission, making each individual's belief vector longer and more resistant
to alternative ideas. Challenging hegemony requires not just presenting
alternative ideas (novelty component) but first \emph{weakening} the
receiver's prior (reducing $w(r)$) --- a process Gramsci called the ``war of
position.''
\end{interpretation}


\section{Projections and Subspace Persuasion}

When we think of persuasion not as a single signal but as a \emph{subspace} of
available signals, the theory becomes even richer.

\begin{definition}[Subspace Persuasion]
\label{def:subspace-persuasion}
A \textbf{persuasion campaign} is a closed subspace $W \subseteq \Hspace$
representing the set of meanings available to the persuader. The
\textbf{optimal persuasive signal} from $W$ for a receiver with prior $r$ is
the element of $W$ that maximizes the persuader's payoff.
\end{definition}

\begin{theorem}[Optimal Signal from a Subspace]
\label{thm:optimal-subspace}
If the persuader can choose any signal with $\embed{s} \in W$ (no norm
constraint), the sender's payoff $\payoff{S}{} = \rs{s}{s} + \rs{s}{r}$ is
unbounded above (just take $\embed{s} = t \cdot P_W(\embed{r}) / \nrm{P_W(\embed{r})}$ for $t \to \infty$).

With a norm constraint $\nrm{\embed{s}} \leq M$, the optimal signal is:
\[
  \embed{s^*} = M \cdot \frac{P_W(\embed{r})}{\nrm{P_W(\embed{r})}},
\]
where $P_W$ is the orthogonal projection onto $W$. That is, the optimal signal
points in the direction of the projection of the receiver's prior onto the
persuader's subspace.
\end{theorem}

\begin{proof}
With norm constraint $M$, the sender's payoff is $M^2 + \ip{\embed{s}}{\embed{r}}$
where $\embed{s} \in W$ with $\nrm{\embed{s}} = M$. We maximize
$\ip{\embed{s}}{\embed{r}}$ over $\embed{s} \in W$ with $\nrm{\embed{s}} = M$.

By Cauchy--Schwarz restricted to $W$: $\ip{\embed{s}}{\embed{r}} =
\ip{\embed{s}}{P_W(\embed{r}) + P_{W^\perp}(\embed{r})} =
\ip{\embed{s}}{P_W(\embed{r})}$ (since $\embed{s} \perp W^\perp$). This is
maximized when $\embed{s} = M \cdot P_W(\embed{r}) / \nrm{P_W(\embed{r})}$.

In Lean~4 (sketch):
\begin{lstlisting}
theorem optimal_signal_subspace (W : Submodule ℝ H)
    (hM : ‖s‖ = M) (hs : s ∈ W) :
    inner s (φ r) ≤ M * ‖orthogonalProjection W (φ r)‖ := by
  sorry -- follows from Cauchy-Schwarz restricted to W
\end{lstlisting}
\end{proof}

\begin{interpretation}[Persuasion as Selective Emphasis]
\label{interp:selective}
The optimal subspace signal theorem has a profound implication for the practice
of persuasion: \emph{the most effective persuasion targets the aspects of the
receiver's beliefs that overlap with the persuader's available messages}.

The persuader does not try to maximize the absolute resonance of their signal
with the receiver. Instead, they identify the component of the receiver's
belief that lies within their own ``message space'' $W$, and they amplify that
component. This is the geometry behind ``selective emphasis'' in rhetoric:
effective persuaders don't fabricate new common ground; they \emph{identify
and amplify existing common ground}.

In political communication, this model explains the practice of ``targeted
messaging.'' A political campaign with a limited set of messages (the subspace
$W$) identifies which of its messages resonate most with each voter segment
(computing $P_W(\embed{r})$ for different $r$) and emphasizes those messages
for those segments. The geometry guarantees that this strategy is optimal.

The model also reveals the fundamental limitation of subspace persuasion: if the
receiver's prior has no component in $W$ (i.e., $P_W(\embed{r}) = 0$, meaning
the receiver's beliefs are entirely orthogonal to the persuader's message
space), then no signal from $W$ can produce positive cross-resonance. The
persuader is ``speaking a different language'' and cannot be effective. This is
the geometric counterpart of the observation that highly specialized propaganda
(say, economic arguments) is ineffective on people whose beliefs are organized
around entirely different dimensions (say, cultural identity).
\end{interpretation}


\section{Iterated Persuasion and Convergence}

What happens when persuasion is applied repeatedly?

\begin{definition}[Iterated Persuasion]
\label{def:iterated}
In \textbf{iterated persuasion}, the sender repeatedly sends signal $s$ to the
receiver, who updates their belief each round:
\begin{align*}
  r_0 &= r \quad \text{(initial prior)}, \\
  r_{n+1} &= r_n \compose s \quad \text{(update)}.
\end{align*}
In the Hilbert space: $\embed{r_n} = \embed{r} + n \embed{s}$.
\end{definition}

\begin{theorem}[Direction Convergence Under Iterated Persuasion]
\label{thm:iterated-convergence}
If $\embed{s} \neq 0$, the direction of the receiver's belief converges to the
direction of the signal:
\[
  \lim_{n \to \infty} \hat{r}_n = \lim_{n \to \infty}
  \frac{\embed{r} + n\embed{s}}{\nrm{\embed{r} + n\embed{s}}}
  = \frac{\embed{s}}{\nrm{\embed{s}}} = \hat{s}.
\]
The convergence rate is $O(1/n)$.
\end{theorem}

\begin{proof}
\[
  \hat{r}_n = \frac{\embed{r} + n\embed{s}}{\nrm{\embed{r} + n\embed{s}}}
  = \frac{\embed{r}/n + \embed{s}}{\nrm{\embed{r}/n + \embed{s}}}
  \xrightarrow{n \to \infty} \frac{\embed{s}}{\nrm{\embed{s}}}.
\]

In Lean~4:
\begin{lstlisting}
theorem iterated_persuasion_conv (hs : φ s ≠ 0) :
    Filter.Tendsto (fun n => (φ r + n • φ s) / ‖φ r + n • φ s‖)
      Filter.atTop (nhds (φ s / ‖φ s‖)) := by
  sorry -- standard analysis argument
\end{lstlisting}
\end{proof}

\begin{interpretation}[The Inevitability of Propaganda]
\label{interp:propaganda}
Theorem~\ref{thm:iterated-convergence} is a mathematical model of propaganda:
\emph{repeated exposure to the same message eventually aligns the receiver's
beliefs with the message, regardless of the receiver's initial beliefs}.

The convergence is unconditional --- it holds for \emph{any} initial prior $r$
and \emph{any} signal $s$ (as long as $s$ is non-void). The receiver's initial
beliefs are irrelevant in the limit; only the signal direction matters. This is
a formalization of the propaganda principle attributed to Joseph Goebbels:
``A lie told once remains a lie, but a lie told a thousand times becomes the
truth.''

But the convergence rate is $O(1/n)$ --- slow, polynomial, not exponential.
A receiver with a strong prior ($w(r) \gg w(s)$) takes many rounds to be
persuaded. This quantifies the ``resistance'' of strong beliefs: you can always
be persuaded eventually, but the stronger your initial conviction, the longer it
takes.

The model also reveals the strategic importance of being the \emph{only}
signal source. If the receiver is exposed to multiple signals from different
directions, the convergence may not occur (the signals may cancel each other
out). Effective propaganda requires not just repetition of one's own message
but \emph{suppression of competing messages} --- a geometric prediction that
accords with the historical practice of censorship in authoritarian regimes.
\end{interpretation}

\begin{theorem}[Competing Persuaders]
\label{thm:competing}
If the receiver is simultaneously exposed to signals $s_1$ and $s_2$ each
round (updating as $r_{n+1} = r_n \compose s_1 \compose s_2$), then:
\[
  \embed{r_n} = \embed{r} + n(\embed{s_1} + \embed{s_2}),
\]
and the receiver's belief converges to the direction of $\embed{s_1} +
\embed{s_2}$ --- the \emph{resultant} of the two signals.
\end{theorem}

\begin{proof}
$\embed{r_n} = \embed{r} + n\embed{s_1 \compose s_2} = \embed{r} +
n(\embed{s_1} + \embed{s_2})$. By the same argument as
Theorem~\ref{thm:iterated-convergence}, $\hat{r}_n \to (\embed{s_1} +
\embed{s_2}) / \nrm{\embed{s_1} + \embed{s_2}}$.
\end{proof}

\begin{interpretation}[The Marketplace of Ideas]
\label{interp:marketplace}
When multiple persuaders compete, the receiver's belief converges to the
\emph{vector sum} of the competing signals. This is a mathematical model of
the ``marketplace of ideas'': in a free society with multiple voices, public
opinion converges to the resultant of all the voices.

The resultant favors voices with larger norms (louder voices have more
influence on the direction) and is sensitive to the angles between voices
(reinforcing voices produce a larger resultant than conflicting ones). This
is a more nuanced picture than the naive ``marketplace'' metaphor suggests:
the outcome is not simply the ``best idea winning'' but a geometric resultant
that depends on the relative weights and directions of all participants.

If the two signals are perfectly opposed ($\embed{s_1} = -\embed{s_2}$), their
resultant is the zero vector: the competing signals cancel, and the receiver's
belief does not converge at all. Their direction remains at $\hat{r}$ forever,
undisturbed by the balanced opposition. This is a model of how balanced media
coverage can paradoxically leave public opinion \emph{unmoved}: when two
equally weighted opposing messages cancel, the net effect is zero.

If one signal is much larger than the other ($\nrm{\embed{s_1}} \gg
\nrm{\embed{s_2}}$), the resultant is approximately in the direction of $s_1$:
the larger signal dominates. This is the geometric foundation of \emph{soft
power}: in a marketplace of ideas, the voice with the most weight (best
production values, widest distribution, most prestigious source) dominates not
by suppressing other voices but by outweighing them vectorially.
\end{interpretation}


\section{Gramsci's Hegemony in Hilbert Space}

We close this chapter with a formalization of Gramsci's concept of cultural
hegemony.

\begin{definition}[Cultural Hegemony]
\label{def:hegemony}
An idea $h \in \IS$ is \textbf{hegemonic} over a population of ideas
$\{a_1, \ldots, a_n\}$ if:
\[
  \rs{h}{a_i} > 0 \quad \text{for all } i,
\]
and
\[
  \nrm{\embed{h}} \geq \nrm{\embed{a_i}} \quad \text{for all } i.
\]
That is, $h$ has positive resonance with every idea in the population and has
at least as much weight as any of them.
\end{definition}

\begin{theorem}[Hegemonic Stability]
\label{thm:hegemonic-stability}
If $h$ is hegemonic over $\{a_1, \ldots, a_n\}$, then for each $i$:
\[
  \payoff{h}{h, a_i} = \rs{h}{h} + \rs{h}{a_i} > \rs{h}{h} > 0.
\]
The hegemonic idea benefits positively from every interaction and benefits more
than any individual counter-idea.
\end{theorem}

\begin{proof}
$\payoff{h}{h, a_i} = \rs{h}{h} + \rs{h}{a_i} > \rs{h}{h}$ since
$\rs{h}{a_i} > 0$. And $\payoff{h}{} - \payoff{a_i}{} = \rs{h}{h} -
\rs{a_i}{a_i} \geq 0$ by the norm dominance condition.
\end{proof}

\begin{interpretation}[The Geometry of Hegemony]
\label{interp:hegemony}
Hegemony in IDT~v2 has a precise geometric meaning: a hegemonic idea is a
vector of large norm that has positive inner product with every idea in the
population. Geometrically, it is a vector that lies in the ``positive orthant''
relative to all population vectors --- it projects favorably onto every cultural
direction simultaneously.

Such vectors exist only when the population's ideas do not span too large
an angular range. If the population includes ideas at angles greater than
$\pi/2$ from each other (partial conflict), no single vector can have positive
inner product with all of them. Hegemony is therefore possible only in
populations with a certain degree of pre-existing ideological coherence ---
a mathematical version of Gramsci's observation that hegemony requires the
``consent'' (here, positive resonance) of the governed.

The norm dominance condition ensures that the hegemonic idea ``outweighs'' every
competing idea, guaranteeing a higher payoff in every pairwise interaction.
This is not merely a matter of truth or persuasive power but of sheer cultural
\emph{weight}: the hegemonic idea is the one that has accumulated the most
norm through historical processes of repetition, institutionalization, and
cultural embedding.

Breaking hegemony requires either (1) introducing an idea that has negative
resonance with $h$ (challenging the hegemony head-on, at the cost of conflict
with the entire population), (2) rotating the population's idea vectors so
that $h$ no longer has positive resonance with all of them (a ``war of
position'' that gradually changes cultural terrain), or (3) introducing a new
dimension of meaning that $h$ does not occupy (an ``outside'' perspective that
creates a new axis of evaluation). All three strategies have historical
precedents, and IDT~v2 can model each geometrically.
\end{interpretation}

\bigskip

Persuasion as projection reveals the deep geometry of influence: reinforcement
vs.\ novelty, the paradox of strong priors, the inevitability of repeated
propaganda, and the conditions for hegemony. In Chapter~\ref{ch:spectral}, we
scale from pairwise interactions to entire networks, using spectral theory to
uncover the principal modes of cultural agreement in populations of ideas.
